%!TEX root=../../main.tex

As in the previous chapters, our analysis will be based on enumerating
the activation patterns associated with ReLU neurons.
However, unlike the previous chapters, we focus on training deep ReLU networks.
As a result, our description of activation patterns and the associated
constraint cones will necessarily be hierarchical.
Before proceeding, we briefly review our notation for deep networks as
introduced in \cref{sec:preliminaries}.

Recall that a deep ReLU network is
parameterized by \( \theta = \rbr{W^{(1)}, \ldots, \Wk} \), where
\( W^{(l)} \in \R^{d_{l} \times d_{l+1}} \) are the weights for the
\( l^{\text{th}} \) layer.
The dimension of the network inputs is given by \( d_1 = d \) and the output
dimension is \( d_{L+1} = c \), where \( c \) is the number of targets.
In accordance with \cref{eq:activations} and \cref{eq:predictions}, the model
activations and predictions are given by
\[
    Z^{(l+1)} = \rbr{\Z W^{(l)}}_+,
    \quad \text{and} \quad
    f_\theta(X) = Z^{(L)} W^{(L)},
\]
respectively, where we use \( Z^{(1)} = X \) for convenience.

Each layer of a neural network is associated with its own set of achievable
activation patterns.
In order to refer to sub-networks,
let \( \theta_l = \rbr{W^{(1)}, \ldots, \Wl} \) be the weights of the
network up to and including the \( l^{\text{th}} \) layer.
The prediction function associated with this sub-network is
\begin{equation}\label{eq:sub-network}
    f^{(l)}_{\theta_{l}}(X) = Z^{(l)} W^{(l)}.
\end{equation}
For a given architecture specified by the hidden dimensions
\( \Delta = (d_2, \ldots, d_{L}) \), we can define the set of activation
patterns at layer \( l \) as follows:
\begin{equation}\label{eq:l-activation-patterns}
    \begin{aligned}
        \calDl_{\Delta} = \big\{
        \Dl_{j_l} =
        \diag\rbr{\mathds{1}\rbr{f^{(l-1)}_{\theta_{l-1}}(X) u \geq 0}}
         & : \theta^{(l-1)} \in \Theta^{(l-1)}, \\
         & u \in \R^{d_{l}}
        \big\}.
    \end{aligned}
\end{equation}
When \( l = 1 \), the sub-network satisfies \( f^{(l)}_{\theta_{l}}(X) = X \)
and \( \calD^{(1)}_{\Delta} \) is fixed (i.e. independent of \( \Delta \))
and equivalent to the standard definition of activation patterns for two-layer
ReLU networks (\cref{eq:activation-patterns}).
For \( l > 1 \), this definition is dependent on \( \Delta \), but taking
all the network widths to infinity yields a canonical set of activations.
\begin{restatable}{lemma}{limitingActivationSet}\label{lemma:activation-patterns}
    The set of activation patterns in the \( l^{\text{th}} \) layer has a
    finite limit, call it \( \calDl \), obtained by taking each layer width
    \( d_{i} \) to be arbitrarily large,
    \[
        \calDl = \lim_{\Delta \rightarrow \infty} \calDl_{\Delta}.
    \]
\end{restatable}
We use \( p_{l} = \abs{\calDl} \) to represent the number of unique activation
patterns in layer \( l \).

Similar to the two-layer setting, we use activation patterns to define
a data-dependent basis expansion on which our convex reformulation will
operate.
Let \( X^{(1)} = X \in \R^{n \times d_1} \) and recursively define
the tensor \( X^{(l+1)} \in \R^{n \times p_1 \times \cdots \times p_{l} \times d_1} \)
such that
\begin{equation}\label{eq:data-tensor}
    X^{(l+1)}_{j_1,\ldots j_l}
    := \Dl_{j_l} \Xl_{j_1, \ldots, j_{l-1}} \in \R^{n \times d_1},
\end{equation}
where we use the convention that \( j_{l} \) always indexes the
\( l^{\text{th}} \) slice of \( X^{(l)} \).
Unrolling this recursion implies that \( X^{(l)} \) is a \( l+1 \) order tensor
satisfying
\( \Xl_{j_1, \ldots, j_{l-1}} = D^{(l-1)}_{j_{l-1}} \cdots D^{(1)}_{j_1} X \).
The order of the activation patterns is not important since activation
patterns are diagonal and so commutative.
Note that for clarity we shall use \( j_l \) as indices of activation patterns
and \( i_l \) as indices of neurons.

The main action of \( \Xl \) will occur in the following tensor reduction
operator.
\begin{definition}[Tensor Reduction]
    Given two tensors
    \( \Xll \in \R^{n \times p_1 \times \cdots \times p_l \times d_1} \)
    and
    \( \Tll \in \R^{d_1 \times p_1 \times \cdots \times d_{l+1}} \),
    the tensor reduction operator denoted by \( \odot \) is defined as
    \begin{equation}\label{eq:tensor-reduction}
        \Xll \odot \Tll = \sum_{j_{l}} \cdots \sum_{j_{1}}
        \Xll_{j_1, \ldots, j_{l}} \Tll_{j_1, \ldots, j_{l}}
    \end{equation}
\end{definition}

Sometimes we will use the tensor reduction with respect to the \( i_{l+1} \)
slice of \( \Tll \), which is given by the following reduction over
matrix-vector products:
\[
    \Xll \odot \Tll_{i_{l+1}} = \sum_{j_{l}} \cdots \sum_{j_{1}}
    \Xll_{j_1, \ldots, j_{l}} \Tll_{j_1, \ldots, j_{l}, i_{l+1}}.
\]
This is particularly useful for defining the tensor cones associated with
each activation pattern \( \Dl_{j_l} \),
\begin{equation}\label{eq:tensor-cone}
    \calK^{(l)}_{j_{l}} =
    \cbr{
    U^{(l)} \in \R^{d_1 \times p_{1} \times \cdots \times p_{l-1}}
    :
    (2D^{(l)}_{j_{l}} - I)
    \sbr{X^{(l)} \odot T^{(l)}_{i_{l}}} \geq 0
    }.
\end{equation}
These cones are a straightforward generalization of \( \calK_j \) from
the two-layer setting seen in Chapters~\ref{chapter:fast-optimization}
and~\ref{chapter:solution-functions}.
In fact, choosing \( l = 1 \) and observing that \( X^{(1)} = X \) and
\( U^{(l)} \in \R^{d_1} \) shows that \( \calK^{(1)}_{j_1} = \calK_j \)
holds exactly.



